% Generated draft: requires mathematical and source review.
\documentclass[11pt]{article}
\usepackage[margin=25mm]{geometry}
\usepackage{fontspec}
\setmainfont{Latin Modern Roman}
\usepackage{amsmath,amssymb,mathtools,unicode-math}
\setmathfont{Latin Modern Math}
\usepackage{enumitem,needspace,xurl,hyperref}
\setlength{\parindent}{0pt}
\setlength{\parskip}{5pt}
\setlength{\emergencystretch}{3em}
\newcommand{\N}{\mathbb{N}}
\newcommand{\Z}{\mathbb{Z}}
\newcommand{\Q}{\mathbb{Q}}
\newcommand{\R}{\mathbb{R}}
\newcommand{\C}{\mathbb{C}}
\newcommand{\F}{\mathbb{F}}
\newcommand{\A}{\mathbb{A}}
\newcommand{\PP}{\mathbb{P}}
\newcommand{\E}{\mathbb{E}}
\newcommand{\eps}{\varepsilon}
\newcommand{\dd}{\,\mathrm{d}}
\begin{document}
\section*{UnsolvedMath ID: 800068}
\textbf{Source problem code:} catalogue entry
\subsection{Definitions and mathematical statement}
An Engel structure on a smooth four-dimensional manifold \(M\) is a smooth rank-two subbundle \(\mathcal D\subset TM\) such that
\[
\mathcal E:=\mathcal D+[\mathcal D,\mathcal D],\qquad
\operatorname{rank}(\mathcal E)=3,\qquad
\mathcal E+[\mathcal D,\mathcal E]=TM.
\]
Here brackets of distributions denote the pointwise spans of Lie brackets of their local smooth sections. Define the full Engel automorphism group by
\[
\operatorname{Aut}(M,\mathcal D)
:=\{f\in\operatorname{Diff}^{\infty}(M):
(df)_p(\mathcal D_p)=\mathcal D_{f(p)}\text{ for every }p\in M\}.
\]
An automorphism is not required to preserve any additional orientation or coorientation.

\textbf{Question.} Does there exist a nonempty closed smooth four-dimensional manifold \(M\), equipped with an Engel structure \(\mathcal D\), such that
\[
\operatorname{Aut}(M,\mathcal D)=\{\operatorname{id}_M\}?
\]
Here closed means compact and without boundary. The smooth category is the conventional interpretation used in this draft; the supplied question does not explicitly specify regularity, connectedness, or orientability, and no connectedness or orientability assumption is imposed here.

\textbf{Supplied context.} The source describes the automorphism group of the Engel local model as isomorphic to the contactomorphism group of standard tight \(\mathbb R^3\), while emphasizing that these automorphisms cannot in general be cut off and globalized. The precise local model and automorphism conventions underlying this contextual identification are not specified in the supplied record. The source attributes examples with very small automorphism groups, and the question of existence of a trivial automorphism group, to R. Montgomery. It reports that Koji Yamazaki recently constructed an example with trivial automorphism group on a noncompact ambient manifold. These are supplied contextual claims, without individual supporting citations in the record. The registry labels the closed-manifold question open; that label is not an independently verified literature assessment.
\subsection{Short English statement}
Does some nonempty compact smooth four-manifold without boundary admit an Engel distribution whose only preserving diffeomorphism is the identity? The supplied source reports a noncompact example and labels the closed case open. This draft uses the full smooth automorphism group and requires human mathematical review.
\subsection{Sources}
\begin{itemize}
\item UnsolvedMath local registry, permanent ID \(800068\), source problem code \texttt{catalogue entry}, title \emph{Foliations and Diffeomorphism Groups --- Question 9}. The supplied record attributes the problem to Nariman and Yazdi, with proposed year 2025, and records status \texttt{open}.
\item Nariman and Yazdi, \emph{Problem List on Foliations and Diffeomorphism Groups} (2025), Question 9. Supplied source URL: \url{https://arxiv.org/abs/2503.18274}. The record reports access on 2026-10-16 and extraction by \texttt{arxiv\_latex}; these are supplied provenance statements, not access or verification performed for this draft.
\item The supplied statement attributes the earlier automorphism-group work and existence question to R. Montgomery, and a noncompact example with trivial automorphism group to Koji Yamazaki. No separate bibliographic references for those works are supplied.
\end{itemize}
\end{document}
